\section*{Chapter \ref{ch:b3:pericyclic} --- Pericyclic Reactions}

\begin{solution}{exo:b3:pericyclic:1}
Diels--Alder: [4+2] \omterm{def:b3:pericyclic:families}{cycloaddition}. Cyclobutene opening: electrocyclic. Cope: [3,3]
sigmatropic. Sulfolene losing $\ce{SO2}$: cheletropic. Ozone and alkene: 1,3-dipolar
([3+2]) \omterm{def:b3:pericyclic:families}{cycloaddition}. $[1,5]$-H shift: sigmatropic.
\end{solution}

\begin{solution}{exo:b3:pericyclic:2}
Hexatriene (6 electrons): thermal \omterm{def:b3:pericyclic:rotation-modes}{disrotatory}, photochemical \omterm{def:b3:pericyclic:rotation-modes}{conrotatory}.
Butadiene (4): thermal \omterm{def:b3:pericyclic:rotation-modes}{conrotatory}, photochemical \omterm{def:b3:pericyclic:rotation-modes}{disrotatory}.
\end{solution}

\begin{solution}{exo:b3:pericyclic:3}
Four electrons, thermal: \omterm{def:b3:pericyclic:rotation-modes}{conrotatory} opening. The two methyl groups, on opposite
faces of the ring, both turn outwards: (\textit{E},\textit{E})-hexa-2,4-diene.
\end{solution}

\begin{solution}{exo:b3:pericyclic:4}
The \omterm{def:b3:pericyclic:faciality}{suprafacial} [2+2] reaction has $4q$ electrons: in the ground state the HOMO of one
ethene and the LUMO of the other overlap with one bonding and one antibonding end.
In the excited state the singly occupied $\pi^*$ of one molecule plays the part of the
HOMO, and both ends match.
\end{solution}

\begin{solution}{exo:b3:pericyclic:5}
The \omterm{def:b3:point-groups:mirror}{mirror plane} is kept. Hexatriene: $\psi_1$ S, $\psi_2$ A, $\psi_3$ S occupied; $\psi_4$ A, $\psi_5$ S, $\psi_6$ A empty.
Cyclohexadiene, in order: $\sigma$ S, $\pi_1$ S, $\pi_2$ A occupied; $\pi_3^*$ S, $\pi_4^*$ A, $\sigma^*$ A empty. Joining by
symmetry in order, $\psi_1 \to \sigma$, $\psi_2 \to \pi_2$, $\psi_3 \to \pi_1$: occupied to occupied, the \omterm{def:b3:pericyclic:rotation-modes}{disrotatory} closure
is thermally allowed.
\end{solution}

\begin{solution}{exo:b3:pericyclic:6}
Six electrons under light: \omterm{def:b3:pericyclic:rotation-modes}{conrotatory}, so the two terminal methyl groups end on
opposite faces: \textit{trans}-5,6-dimethylcyclohexa-1,3-diene (the thermal reaction gives
the \textit{cis} isomer).
\end{solution}

\begin{solution}{exo:b3:pericyclic:7}
The hydrogen on C5 moves to C1 of the diene, its neighbour around the ring,
through a six-electron transition state ($\sigma2_s + \pi4_s$), \omterm{def:b3:pericyclic:faciality}{suprafacially}: one migration
gives 1-methyl-, a second 2-methylcyclopentadiene. A five-carbon chain allows the
hydrogen to stay on one face, which the small ring forces anyway.
\end{solution}

\begin{solution}{exo:b3:pericyclic:8}
$\pi2_s$ (alkene) $+ \pi2_a$ (ketene C=C, attacked on opposite faces through the
perpendicular C=O $\pi^*$): the count of $(4q + 2)_s$ components is one (the alkene), odd:
thermally allowed.
\end{solution}

\begin{solution}{exo:b3:pericyclic:9}
In the chair, the C=C of the vinyl group (atoms 1--2), the oxygen (3), the $\ce{OCH2}$
carbon (4) and the allyl C=C (5--6); the O--C bond breaks and the C1--C6 bond forms:
the product is pent-4-enal,
$\ce{OHC-CH2-CH2-CH=CH2}$.
\end{solution}

\begin{solution}{exo:b3:pericyclic:10}
Eight electrons is $4q$; with one \omterm{def:b3:pericyclic:faciality}{antarafacial} component the array is of Möbius
topology, aromatic with $4q$ electrons: the transition state is aromatic and the
reaction thermally allowed (as the \omterm{def:b3:pericyclic:rotation-modes}{conrotatory} closure of octatetraene).
\end{solution}

\begin{solution}{exo:b3:pericyclic:11}
Pairing the largest coefficients joins N3 to the terminal carbon: the
1,4-disubstituted triazole. Without a catalyst the coefficients differ little and
both 1,4 and 1,5 isomers form; copper(I) makes the reaction stepwise through a
copper acetylide and gives the 1,4 isomer only.
\end{solution}

\begin{solution}{exo:b3:pericyclic:12}
The hydrogen bridges the termini of the SOMO of a $j$-atom radical, $k = (j + 1)/2$. Its terminal
coefficients are proportional to $\sin(k\pi/(j + 1))$ and $(-1)^{k+1}\sin(k\pi/(j + 1))$. For $j = 7$, $k = 4$: opposite signs, so
the lobes of equal sign are on opposite faces at the two ends: \omterm{def:b3:pericyclic:faciality}{antarafacial}. For
$j = 5$, $k = 3$: the same sign on the same face: \omterm{def:b3:pericyclic:faciality}{suprafacial}.
\end{solution}

\begin{solution}{pb:b3:pericyclic:1}
\textbf{1.} The C9--C10 $\sigma$ bond of ring B; with the 5,7-diene it gives a hexatriene,
C10=C5--C6=C7--C8=C9.

\textbf{2.} Six electrons (the two $\pi$ bonds and the $\sigma$ bond): an electrocyclic ring opening.

\textbf{3.} Thermally, a cyclohexadiene and its open hexatriene equilibrate only
slowly, through a high barrier, and the ring is the more stable side; the
diene absorbs ultraviolet light, and its excited state opens within
picoseconds.

\textbf{4.} \omterm{def:b3:pericyclic:rotation-modes}{Conrotatory}: in the excited state the singly occupied LUMO of the triene
system ($k = 4$) has termini of opposite signs.

\textbf{5.} The C6=C7 bond comes from the ring: its two chain continuations, C5 and C8,
were on the same side of it, and stay so: \textit{Z}.

\textbf{6.} The thermal opening, allowed only disrotatorily, has a barrier far beyond what
body heat can supply, and would lead uphill, towards the less stable triene.

\textbf{7.} From C19 (the methyl on C10) to C9: a $[1,7]$-H sigmatropic shift.

\textbf{8.} Eight: the six $\pi$ electrons of the triene and the two of the C--H bond.

\textbf{9.} $\sigma2 + \pi6$. \omterm{def:b3:pericyclic:faciality}{Suprafacially} both would be s, with one $(4q + 2)_s$ component from each:
even, forbidden. With the triene \omterm{def:b3:pericyclic:faciality}{antarafacial}, $\sigma2_s + \pi6_a$: one counted component, odd,
allowed.

\textbf{10.} Seven atoms in a helical, \textit{Z}-configured chain let the hydrogen pass from one face to
the other; a three-atom chain cannot twist that far.

\textbf{11.} Möbius (one \omterm{def:b3:pericyclic:faciality}{antarafacial} component) with eight electrons, $4q$: aromatic.

\textbf{12.} It is thermally allowed; its barrier is crossed slowly at body temperature.

\textbf{13.} Six electrons under light: \omterm{def:b3:pericyclic:rotation-modes}{conrotatory} again, but it can turn the termini the
other way and put H9 and the C19 methyl on the other faces: a ring-closed
stereoisomer, lumisterol.

\textbf{14.} With the central bond \textit{E}, C19 and C9 are on opposite sides of the chain and far
apart: the hydrogen cannot reach.

\textbf{15.} Previtamin D$_3$ itself absorbs and is converted into lumisterol and tachysterol,
which do not give vitamin D; a photostationary mixture is reached, which caps the
production.

\textbf{16.} The opening happens in the excited state, in femtoseconds to picoseconds;
the shift is a thermal reaction with a barrier of \qty{85}{kJ/mol}.

\textbf{17.} $t_{1/2} = \ln 2/k = 0.693/\qty{0.023}{h^{-1}} = \qty{30}{h}$.

\textbf{18.} $1 - \eu^{-0.023 \times 24} = 0.42$.

\textbf{19.} $\ln 10/k = 2.303/\qty{0.023}{h^{-1}} = \qty{100}{h}$.

\textbf{20.} $k(\qty{293.15}{K}) = k(\qty{310.15}{K})\,\eu^{-(E_a/R)(1/293.15 - 1/310.15)} = 0.023 \times \eu^{-1.91} = \qty{3.4e-3}{h^{-1}}$.

\textbf{21.} $2.303/\qty{3.4e-3}{h^{-1}} = \qty{680}{h}$, about 28 days.

\textbf{22.} The thermal step is seven times faster at \qty{37}{\celsius} than at \qty{20}{\celsius}, so the
previtamin made in a morning in the sun is converted within days.

\textbf{23.} The steroid skeleton, with its many stereocentres, comes ready-made from
the sterol; light and warmth then carry out the two pericyclic steps exactly as in
the skin.

\textbf{24.} At body temperature, 90\,\% of the previtamin D$_3$ becomes vitamin D$_3$ in about
\qty{100}{h}, four days.
\end{solution}
